\section{To'yingan bug'}


\begin{enumerate}

\item Suyuqlikdan biror vaqt ichida chiqib ketayotgan molekulalar soni o‘sha vaqt ichida suyuqlikka qaytib tushayotgan bug‘ molekulalari soniga teng bo‘lishiga nima deyiladi?
\begin{enumerate}[label=\Alph*)]
\item Termodinamik muvozanat
\item Kyuri holati
\item Dinamik muvozanat
\item Kritik holat
\end{enumerate}

\item To‘yingan bug‘ bosimining hajmga bog‘lanishi qanday?
\begin{enumerate}[label=\Alph*)]
\item hajmga to‘g‘ri proporsional
\item hajmga bog‘liq emas
\item hajmga teskari proporsional
\item hajmning kvadratiga proporsional
\end{enumerate}

\item Germetik berk idishda faqat to‘yingan suv bug‘i bor (suv yo‘q). Idish isitilganda bug‘ molekulalarining konsentratsiyasi qanday o‘zgaradi?
\begin{enumerate}[label=\Alph*)]
\item o‘zgarmaydi
\item kamayadi
\item ortadi
\item ortishi ham kamayishi ham mumkin
\end{enumerate}

\item 1-yopiq idishda suv va suv bug‘i bor, ikkinchi idishda faqat to‘yingan suv bug‘i bor. Harorat ortganida bu idishlardagi bosim qanday o‘zgaradi?
\begin{enumerate}[label=\Alph*)]
\item Ikkalasida bir xil ortadi
\item 1-sida ko‘proq ortadi
\item 2-sida ko‘proq ortadi
\item 1-sida o‘zgarmaydi, 2-sida ortadi
\end{enumerate}

\item Moddaga tegishli bo‘lgan kritik temperaturadan yuqori haroratda modda qanday agregat holatida bo‘ladi?
\begin{enumerate}[label=\Alph*)]
\item gaz
\item suyuqlik
\item qattiq
\item to‘yingan bug‘
\end{enumerate}

\item Qaysi holda gazlarni suyuqlikka aylantirish mumkin bo‘ladi?
\begin{enumerate}[label=\Alph*)]
\item har qanday temperaturada
\item kritik temperaturadan yuqori temperaturada
\item faqat $0^\circ C$ haroratda
\item kritik temperaturadan past temperaturada
\end{enumerate}

\item Berilgan jumlalarning qaysi biri noto‘g‘ri?
\begin{enumerate}[label=\Alph*)]
\item to‘yingan bug‘ Mendeleyev-Klapeyron qonuniga bo‘ysunadi.
\item kritik temperaturada suyuqlik bilan to‘yingan bug‘ning fizik xossalari farqi yo‘qoladi.
\item temperatura qancha yuqori bo‘lsa, to‘yingan bug‘ning elastikligi shuncha katta bo‘ladi.
\item to‘yingan bug‘ning elastikligi o‘zgarmas temperaturada bug‘ egallab turgan hajmga bog‘liq.
\end{enumerate}



\item O‘zgarmas temperaturada to‘yingan bug‘ning hajmi 4 marta kamaytirilsa, bosim qanday o‘zgaradi?
\begin{enumerate}[label=\Alph*)]
\item 2 marta oshadi
\item 2 marta kamayadi
\item o‘zgarmaydi
\item 4 marta kamayadi
\end{enumerate}

\item Yopiq idishda suv va uning to‘yingan bug‘i bor. Harorat pasayganda to‘yingan bug‘ning zichligi qanday o‘zgaradi?
\begin{enumerate}[label=\Alph*)]
\item kamayadi
\item ortadi
\item o‘zgarmaydi
\item ortishi ham kamayishi ham mumkin
\end{enumerate}

\item To‘yingan suv bug‘ining temperaturasi $100^\circ C$, u biror hajmni egallab turibdi. Dastlabki temperaturani saqlagan holda bug‘ning hajmini ikki marta kamaytirsak, uning bosimi qanday o‘zgaradi?
\begin{enumerate}[label=\Alph*)]
\item o‘zgarmaydi
\item ortadi
\item kamayadi
\item 2 marta ortadi
\end{enumerate}

\item Idishda suv va uning to‘yingan bug‘i bor. Harorat o‘zgartirilmasdan to‘yingan bug‘ning hajmi ikki marta kamaytirilsa, uning bosimi qanday o‘zgaradi?
\begin{enumerate}[label=\Alph*)]
\item 1,4 marta kamayadi
\item 1,4 marta ortadi
\item o‘zgarmaydi
\item 2 marta ortadi
\end{enumerate}

\item Kritik temperaturada suyuqlikning zichligi $925 \, kg/m^3$. Shu temperaturada uning bug‘ining zichligini toping ($kg/m^3$).
\begin{enumerate}[label=\Alph*)]
\item 1
\item 9,25
\item 92,5
\item 925
\end{enumerate}

\item $100^\circ C$ da to‘yingan suv bug‘ining zichligi qanday bo‘ladi ($kg/m^3$)?
\begin{enumerate}[label=\Alph*)]
\item 0,59
\item 0,28
\item 0,16
\item 0,75
\end{enumerate}

\item Kritik temperaturada suyuqlikning zichligi $850 \, kg/m^3$. Shu temperaturada uning bug‘ining zichligini toping ($kg/m^3$).
\begin{enumerate}[label=\Alph*)]
\item 1
\item 0,825
\item 8,25
\item 850
\end{enumerate}

\item To‘yingan bug‘ning bosimi 29,5 kPa. Uning hajmini ikki marta orttirsak 2 soatdan keyin bosimi nimaga teng bo‘ladi (kPa). $T = const$.
\begin{enumerate}[label=\Alph*)]
\item 1
\item 100
\item 14,75
\item 29,5
\end{enumerate}

\end{enumerate}



\begin{table}[h!]
\centering
\begin{tabular}{|c|c|c|c|c|c|c|c|c|c|}
\hline
№ & J. & № & J. & № & J. & № & J. & № & J. \\ \hline
1 & C & 5 & A & 9 & A & 13 & A & & \\ \hline
2 & B & 6 & D & 10 & A & 14 & D & & \\ \hline
3 & A & 7 & D & 11 & C & 15 & D & & \\ \hline
4 & B & 8 & C & 12 & D &  &  & & \\ \hline
\end{tabular}
\end{table}



\section{Havoning namligi. Nisbiiy namlik}


\begin{enumerate}

\item Ta’rifni davom ettiring: "Havodagi suv bug‘ining parsial bosimi bu........."
\begin{enumerate}[label=\Alph*)]
\item bug‘ kondensatsiyalanadigan bosim
\item suv bug‘i to‘yinganda hosil qiladigan bosim
\item boshqa gazlar bo‘lmaganda suv bug‘i hosil qiladigan bosim
\item havoning barometr ko‘rsatadigan bosimi
\end{enumerate}

\item Gapni davom ettiring. Absolyut namlik deb:
1) berilgan temperaturadagi to‘yingan bug‘ bosimiga aytiladi.
2) berilgan temperaturada havodagi suv bug‘ining parsial bosimiga aytiladi.
3) havodagi to‘yingan bug‘ massasiga aytiladi.
4) $1 m^3$ havodagi suv bug‘larining miqdoriga aytiladi.
5) havoda bo‘lgan suv bug‘ining miqdori bilan o‘lchanadigan kattalikka aytiladi.
\begin{enumerate}[label=\Alph*)]
\item 1, 5
\item 2, 4, 5
\item 2, 3
\item 3, 4
\end{enumerate}

\item Shudring nuqtasi qaysi asbob yordamida aniqlanadi?
\begin{enumerate}[label=\Alph*)]
\item ariometr
\item manometr
\item taxometr
\item gigrometr
\end{enumerate}

\item Quyida keltirilgan fikrlarning qaysi biri noto‘g‘ri?
1) shudring nuqtasida havoning nisbiy namligi 100\% oshadi,
2) suv bug‘i toyinadigan holdagi temperatura shudring nuqtasi deb ataladi
3) havo shudring nuqtasigacha sovisa bug‘ kondensatsiyalanib boshlaydi: tuman hosil bo‘ladi, shudring tushadi.
4) shudring nuqtasida havodagi suv bug‘larining elastligi to‘yingan bug‘ elastligiga teng.
5) shudring nuqtasida havodagi suv bug‘larining zichligi to‘yingan bug‘ zichligiga teng.
\begin{enumerate}[label=\Alph*)]
\item 1
\item 2
\item 3
\item 4
\end{enumerate}

\item Ichida namlikni yutuvchi modda bo‘lgan naycha orqali 20 $l$ havo o‘tkazilgan. Bunda naychaning massasi 400 mg ortdi. Havoning absolyut namligi qanday ($g/m^3$)?
\begin{enumerate}[label=\Alph*)]
\item 40
\item 30
\item 20
\item 50
\end{enumerate}

\item Ichida namlikni yutuvchi modda bo‘lgan naycha orqali 2 $l$ havo o‘tkazilganda havoning absolyut namligi $60 \, g/m^3$ ekanligi aniqlandi. Bunda naychaning massasi qanchaga ortgan?
\begin{enumerate}[label=\Alph*)]
\item 30 g
\item 12 g
\item 12 mg
\item 120 mg
\end{enumerate}

\item $20^\circ C$ temperaturada $5 \, m^3$ havoda 50 g suv bug‘i bo‘lsa, havoning nisbiy namligi qanday? ($20^\circ C$ temperaturada to‘yingan suv bug‘ining zichligi $17,3 \, g/m^3$).
\begin{enumerate}[label=\Alph*)]
\item 65\%
\item 62\%
\item 58\%
\item 76\%
\end{enumerate}

\item Namligi 15\% bo‘lgan $1 \, m^3$ havo namligi 30\% bo‘lgan $2 \, m^3$ havo bilan aralashtirildi. Aralashma o‘sha haroratda $3 \, m^3$ hajmni egallaydi. Aralashmaning namligini toping.
\begin{enumerate}[label=\Alph*)]
\item 25\%
\item 21,5\%
\item 26,7\%
\item 28\%
\end{enumerate}

\item Sig‘imi 2 litr bo‘lgan berk idishda $20^\circ C$ temperaturada to‘yingan suv bug‘i bor. Temperatura $5^\circ C$ gacha pasayishida idishda qancha suv hosil bo‘ladi? (to‘yingan bug‘ning zichligi $20^\circ C$ da $17,3 \, g/m^3$; $5^\circ C$ da $6,8 \, g/m^3$).
\begin{enumerate}[label=\Alph*)]
\item 13,6 mg
\item 21 mg
\item 10,6 mg
\item 5,33 mg
\end{enumerate}

\item Xona havosining namligi 90\%. Agar psixrometrning $T_1$ – quruq termometrning ko‘rsatkichi, $T_2$ – ho‘l termometrning ko‘rsatkichi bo‘lsa, quyidagi munosabatlarning qaysi biri to‘g‘ri?
\begin{enumerate}[label=\Alph*)]
\item $T_1 > T_2$
\item $T_1 < T_2$
\item $T_1 = T_2$
\item $T_1 = T_2 + 90K$
\end{enumerate}

\item Psixrometrning quruq va nam termometrlari ko‘rsatkichlarining farqi ortsa, havoning nisbiy namligi qanday o‘zgaradi?
\begin{enumerate}[label=\Alph*)]
\item kamayadi
\item ortadi
\item o‘zgarmaydi
\item ortishi ham kamayishi ham mumkin
\end{enumerate}

\item Psixrometrning quruq va nam termometrlari ko‘rsatkichlarining farqi kamaysa, havoning nisbiy namligi qanday o‘zgaradi?
\begin{enumerate}[label=\Alph*)]
\item kamayadi
\item ortadi
\item o‘zgarmaydi
\item ortishi ham kamayishi ham mumkin
\end{enumerate}

\item Psixrometr quruq termometrni temperaturasi o‘zgarmas qoldirib nam termometr temperaturasini pasaytirsak havoning nisbiy namligi qanday o‘zgaradi?
\begin{enumerate}[label=\Alph*)]
\item kamayadi
\item ortadi
\item o‘zgarmaydi
\item ortishi ham kamayishi ham mumkin
\end{enumerate}

\item Havoning absolyut namligini aniqlang ($kg/m^3$). Havodagi bug‘ning parsial bosimi 14 kPa, temperatura $60^\circ C$ (T = 333 K). Suvning molyar massasi 0,018 kg/mol.
\begin{enumerate}[label=\Alph*)]
\item 0,065
\item 0,091
\item 0,65
\item 0,340
\end{enumerate}

\item Xonadagi havoning temperaturasi $19^\circ C$ ga suv bug‘ining partsial bosimi esa $1,1 \, kPa$ ga teng. Agar bu temperaturada to‘yingan suv bug‘ining bosimi $2,2 \, kPa$ ga teng bo‘lsa, shu xonadagi nisbiy namlik qanday \% ?
\begin{enumerate}[label=\Alph*)]
\item 40
\item 50
\item 55
\item 60
\end{enumerate}

\item Sig‘imi 2 $l$ bo‘lgan berk idishda $20^\circ C$ da to‘yingan suv bug‘i bor. Temperatura $5^\circ C$ gacha pasayishida idishda qancha suv hosil bo‘ladi (mg)? 278 K, 291 K, 293 K temperaturada to‘yingan bug‘ zichliklari mos ravishda $6,8 \, g/m^3, 15,4 \, g/m^3, 17,3 \, g/m^3$.
\begin{enumerate}[label=\Alph*)]
\item 12
\item 21
\item 42
\item 82
\end{enumerate}

\item To‘yingan suv bug‘i molekulalarining konsentratsiyasi $10^\circ C$ dagiga qaraganda $20^\circ C$ da necha marta katta bo‘ladi? 283 K, 291 K, 293 K temperaturada to‘yingan bug‘ zichliklari mos ravishda $9,4 \, g/m^3, 15,4 \, g/m^3, 17,3 \, g/m^3$.
\begin{enumerate}[label=\Alph*)]
\item 2,83 marta
\item 1,83 marta
\item 1,21 marta
\item 1,12 marta
\end{enumerate}

\item Kunduzi $20^\circ C$ da havoning nisbiy namligi 60\% bo‘lgan. Agar tunda temperatura $8^\circ C$ gacha pasaysa, har bir kub metr havodan qancha suv shudring ko‘rinishida ajralib chiqadi (g)? 281 K, 284 K, 293 K temperaturada to‘yingan bug‘ zichliklari mos ravishda $8,3 \, g/m^3, 10 \, g/m^3, 17,3 \, g/m^3$.
\begin{enumerate}[label=\Alph*)]
\item 0,4
\item 21
\item 3,4
\item 2,1
\end{enumerate}

\item Kechqurun $16^\circ C$ da havoning nisbiy namligi 55\% ga teng. Agar tunda temperatura $8^\circ C$ gacha pasaysa, shudring tushadimi? 281 K, 284 K, 289 K temperaturada to‘yingan bug‘ zichliklari mos ravishda $8,3 \, g/m^3, 10 \, g/m^3, 13,6 \, g/m^3$.
\begin{enumerate}[label=\Alph*)]
\item Shudring tushmaydi
\item Ha tushadi
\item Tushishi ham tushmasligi ham mumkin
\item Noma'lum
\end{enumerate}

\end{enumerate}



\begin{table}[h!]
\centering
\begin{tabular}{|c|c|c|c|c|c|c|c|c|c|}
\hline
№ & J. & № & J. & № & J. & № & J. & № & J. \\ \hline
1 & C & 5 & C & 9 & B & 13 & A & 17 & B \\ \hline
2 & B & 6 & D & 10 & A & 14 & B & 18 & D \\ \hline
3 & D & 7 & C & 11 & A & 15 & B & 19 & A \\ \hline
4 & A & 8 & A & 12 & B & 16 & B & & \\ \hline
\end{tabular}
\end{table}
